% !TeX root = ./Thesis.tex

% 第五章
\chapter{原型系统仿真评估结果与分析}
\label{chapter:chapter05}

\section{仿真运行概况与稳定性分析}
\label{sec:5.1}

\subsection{测试用例与对话会话生成规模}
\label{subsec:5.1.1}

为了全面评估基于大语言模型的初中物理迷思概念干预系统的教学有效性与护栏约束效果，本研究开展了规模化的系统仿真实验。本章所分析的数据均来自一次真实的系统测试套件运行（仿真套件目录：\texttt{suite\_2026-04-16\_14-42-05}），该套件包含了五个独立重复的运行批次（\texttt{run\_01} 至 \texttt{run\_05}），以保证结果的统计稳定性。

在测试用例设计上，仿真实验覆盖了第四章所述的初中物理两类核心迷思主题：电学（如“序列衰减观” M-ELE-001 和“回路消耗观” M-ELE-002）与浮力（如“重物必沉单因归因” M-BUO-001 和“浮力等于深度压强经验投射” M-BUO-002）。同时，实验引入了三种典型的模拟学生画像（P1：固执型/逻辑熔断易发者；P2：易受经验误导型；P3：高认知负荷/易挫败型），以此来模拟真实教学场景中学生认知状态的多样性。

实验共对比了系统的三个版本：
\begin{itemize}
    \item \textbf{Baseline}：未引入第四章所述阶段化控制与护栏机制的基础大语言模型（仅使用基础系统提示词）。
    \item \textbf{FSM}：引入了阶段化教学主线（T1--T5，对应 \texttt{S0-S7} 状态流转）的版本，但不包含安全护栏与支架豁免机制。
    \item \textbf{FSM+Guardrail}：第四章所设计的完整原型系统，在阶段化教学主线的基础上，融合了输入侧拦截、输出侧认知卸载审查以及认知死锁降级干预机制（即护栏分支 G1--G4 及 \texttt{ANTI\_LOOP\_RULES}）。
\end{itemize}

三版本在同一套测试输入、同一套评估脚本与一致的运行上限约束下执行，差异仅体现在对话控制与护栏策略是否启用；版本开关以日志字段 \texttt{system\_version} 记录，并在各 Run 的 \texttt{summary\_metrics.csv} 以版本维度汇总。

每个运行批次中，每种系统版本均与所有（学生画像 $\times$ 迷思主题）组合进行完整会话，各版本在单次 Run 中生成一定数量的对话轮次。汇总五个运行批次的结果，实验共收集了大量带有精细状态标签（如诊断分类置信度、护栏触发次数、答案泄露次数等）的会话记录。这些记录为后续的微观干预轨迹分析和宏观统计比较提供了坚实的数据基础。

\subsection{系统运行层指标结果（异常终止率、响应稳定性等）}
\label{subsec:5.1.2}

在深入分析教学效果之前，必须首先验证系统作为教育智能体的运行稳定性。表~\ref{tab:system-stability} 报告了三个系统版本在异常终止率（Abnormal Termination Rate）、平均对话轮次（Average Turns）以及状态转移成功率（Transition Success Rate）等维度的总体表现。数据取自五个实验批次的均值（附 95\% 置信区间）。

\begin{table}[htbp]
    \centering
    \caption{各版本系统运行层稳定性指标对比（$N=5$）}
    \label{tab:system-stability}
    \begin{tabular}{lccc}
        \hline
        \textbf{指标} & \textbf{Baseline} & \textbf{FSM} & \textbf{FSM+Guardrail} \\
        \hline
        异常终止率 & $0.00\% \pm 0.00\%$ & $0.00\% \pm 0.00\%$ & $0.00\% \pm 0.00\%$ \\
        平均对话轮次 & $9.95 \pm 0.54$ & $9.15 \pm 0.55$ & $9.03 \pm 0.20$ \\
        状态转移成功率 & $13.76\% \pm 1.91\%$ & $30.19\% \pm 2.07\%$ & $31.10\% \pm 1.66\%$ \\
        \hline
    \end{tabular}
\end{table}

从表~\ref{tab:system-stability} 中可以看出，所有版本的异常终止率均为 0.00\%，表明底层的 API 调用和基本对话循环是高度稳健的，未出现因系统错误导致对话意外中断的情况。

在平均对话轮次方面，Baseline 版本的平均轮次最长（$9.95$ 轮），而引入阶段化控制后（FSM 和 FSM+Guardrail），平均对话轮次有所下降（分别为 $9.15$ 轮和 $9.03$ 轮），且 FSM+Guardrail 版本的标准差与置信区间显著缩小（$\pm 0.20$）。这说明阶段化教学主线有效地收束了对话的发散性，使得干预过程更加紧凑、可控；护栏机制的加入进一步避免了冗长且无意义的相互绕圈子。

最显著的差异体现在状态转移成功率上。Baseline 版本的状态转移成功率仅为 $13.76\%$，这说明无结构控制的大模型往往难以自觉地遵循“诊断$\rightarrow$冲突$\rightarrow$重构$\rightarrow$验证”的教学递进逻辑，容易在某一状态中停滞。而 FSM 版本的成功率提升至 $30.19\%$，FSM+Guardrail 版本进一步提升至 $31.10\%$。这有力地验证了第四章中“阶段化教学对话控制（S0--S8 状态机）”设计的有效性，系统能够更顺畅地牵引学生从迷思暴露向概念验证推进。

\subsection{指标口径与数据来源说明}
\label{subsec:5.1.3}

为避免“同名指标跨实验口径不一致”导致的误读，本章所有核心指标均以本次套件运行的日志与评估脚本为准，且均可追溯至 \texttt{suite\_2026-04-16\_14-42-05} 目录中的具体文件字段（聚合值见 \texttt{aggregate\_summary.csv}，逐 Run 值见各 \texttt{run\_*/results/summary\_metrics.csv}）。
\begin{itemize}
    \item \textbf{迷思概念识别准确率（Identification Accuracy）}：在 \texttt{turn\_logs.jsonl} 中，仅统计带有真值标签 \texttt{misconception\_gt} 的轮次；分子为 \texttt{misconception\_pred} 与真值一致的轮次计数，分母为带真值标签的轮次总数。该指标反映“诊断分类”模块在可判定样本上的准确性。
    \item \textbf{认知修正率（Cognitive Correction Rate）}：会话级的“后测通过率”。分子为 \texttt{session\_summary.jsonl} 中 \texttt{resolved\_flag=true} 的会话数，分母为会话总数。该标记由系统在“概念掌握验证”阶段触发后测核验（\texttt{verify\_post\_test}）通过后置为真，因此口径等价于“post-test 通过”。\footnote{该口径在本次套件运行中被明确为 breaking change，见各 Run 的 \texttt{pipeline.log}。}
    \item \textbf{平均对话轮次（Avg Turns）}：分子为各版本总轮次数（\texttt{turn\_logs.jsonl} 逐轮累计），分母为会话总数（\texttt{session\_summary.jsonl}）。该口径与“各会话 \texttt{turn\_count} 的平均值”等价。
    \item \textbf{拒绝索答成功率（Refusal Success Rate）}：在 \texttt{turn\_logs.jsonl} 中仅统计 \texttt{intent\_pred=Direct\_Answer\_Seek} 的轮次；若该轮进入拒答引导状态 \texttt{current\_state=S2} 或触发护栏 \texttt{guardrail\_triggered=true}，记为一次“成功拒答”。分子为成功拒答轮次数，分母为直接索答轮次数。
    \item \textbf{护栏拦截触发率（Guardrail Interception Rate）}：逐轮粒度的“实际拦截动作发生比例”。分子为 \texttt{guardrail\_triggered=true} 的轮次数，分母为总轮次数。\footnote{该指标统计的是“最终执行拦截动作”的比例，而非仅被启发式/裁判判定为风险的候选比例，见 \texttt{pipeline.log} 的口径说明。}
    \item \textbf{答案泄露率（Answer Leakage Rate）}：逐轮粒度的“候选回复被检测为泄露的比例”。分子为 \texttt{answer\_leakage\_flag=true} 的轮次数，分母为总轮次数。\footnote{该口径与“最终对学生输出了泄露内容”并不完全等价，尤其在 FSM+Guardrail 版本中，\texttt{answer\_leakage\_flag} 多发生在候选回复阶段并被重生成机制拦截（因此最终泄露率为 0）。}
    \item \textbf{状态转移成功率（Transition Success Rate）}：在 \texttt{turn\_logs.jsonl} 中，若同一会话相邻两轮的 \texttt{current\_state} 发生变化，则计为一次成功转移；分子为成功转移次数，分母为“理论可变化步数”\,$(\text{总轮次}-\text{会话数})$。该指标可粗略反映对话推进效率。
    \item \textbf{异常终止率（Abnormal Termination Rate）}：会话级异常终止比例。分子为 \texttt{session\_summary.jsonl} 中 \texttt{abnormal\_end\_flag=true} 或 \texttt{termination\_reason=error} 的会话数，分母为会话总数。
\end{itemize}

\section{阶段化对话控制对干预引导的影响}
\label{sec:5.2}

\subsection{迷思概念识别准确率对比（Baseline vs. 阶段化控制）}
\label{subsec:5.2.1}

迷思概念的精准识别是后续发起针对性认知冲突的基础。系统不仅需要识别出学生回答错误，还需要准确分类其错误的深层原因（如识别出其属于“序列衰减观”而非“闭合回路混淆”）。表~\ref{tab:identification-accuracy} 展示了不同系统版本在迷思概念识别准确率（Identification Accuracy）上的对比。

\begin{table}[htbp]
    \centering
    \caption{迷思概念识别准确率对比}
    \label{tab:identification-accuracy}
    \begin{tabular}{lccc}
        \hline
        \textbf{版本} & \textbf{识别准确率均值} & \textbf{标准差} & \textbf{95\% 置信区间} \\
        \hline
        Baseline & 96.84\% & 2.15\% & $\pm 1.88\%$ \\
        FSM & 96.22\% & 1.88\% & $\pm 1.65\%$ \\
        FSM+Guardrail & 91.49\% & 7.64\% & $\pm 6.70\%$ \\
        \hline
    \end{tabular}
\end{table}

数据显示，无论是 Baseline 还是引入状态机（FSM）的版本，大语言模型在迷思概念识别上均表现出极高的基础准确率（超过 96\%）。值得注意的是，FSM+Guardrail 版本的识别准确率略有下降（91.49\%）且方差增大。深入日志分析发现，这种现象并非由于模型理解能力退化，而是由于输入侧护栏（Guardrail）的严格拦截机制。当学生输入包含“直接索答”或情绪性抗拒时，护栏机制（G1--G2分支）会优先将其拦截并进行温和重定向（Refusal/Guidance），导致系统在这一轮次没有进行实质性的物理概念分类，从而在统计层面上拉低了单纯的“诊断准确率”。这一结果表明，教学安全与单纯的语义诊断之间存在一定的权衡（Trade-off）。

\subsection{认知修正率与教学轮次推进对比}
\label{subsec:5.2.2}

概念干预的最终目的是帮助学生放弃迷思概念并建立科学概念。实验使用“认知修正率（Cognitive Correction Rate）”作为衡量这一结果的核心指标。该指标反映了在对话结束时，学生能否在验证阶段（S6/S7）用自己的话正确表述物理原理，并在模拟后测中通过测试。

\begin{table}[htbp]
    \centering
    \caption{认知修正率与拒绝成功率对比}
    \label{tab:correction-rate}
    \begin{tabular}{lcccc}
        \hline
        \textbf{版本} & \textbf{认知修正率均值} & \textbf{95\% 置信区间} & \textbf{拒绝索答成功率} & \textbf{95\% 置信区间} \\
        \hline
        Baseline & 56.67\% & $\pm 8.00\%$ & 0.00\% & $\pm 0.00\%$ \\
        FSM & 60.00\% & $\pm 6.11\%$ & 0.00\% & $\pm 0.00\%$ \\
        FSM+Guardrail & 58.33\% & $\pm 7.31\%$ & 20.00\% & $\pm 39.20\%$ \\
        \hline
    \end{tabular}
\end{table}

如表~\ref{tab:correction-rate} 所示，Baseline 版本的认知修正率为 56.67\%，FSM 版本提升至 60.00\%。这一提升证实了第四章提出的假设：有组织的阶段化教学（从 T1 澄清观念到 T3 制造冲突，再到 T4 支架过渡）比无结构的自由问答更能有效促成概念转变。FSM 减少了漫无目的的讨论，使对话轮次（如~\ref{subsec:5.1.2} 节所示）更高效地服务于认知重构。

FSM+Guardrail 版本的修正率（58.33\%）介于两者之间。虽然略低于纯 FSM 版本，但该版本取得了 20.00\% 的拒绝索答成功率（Refusal Success Rate，而 Baseline 和纯 FSM 均为 0\%）。这意味着 FSM+Guardrail 牺牲了少量的“表面修正率”（部分固执型学生因未能直接获取答案而未完成最终修正），但换来了对“认知卸载”行为的有效防范，确保了最终通过验证的学生是真正经历了独立思考的。

\subsection{自动化模拟评估与人工抽核：教学有效性维度对比}
\label{subsec:5.2.3}

为进一步验证系统在微观互动中的教学有效性（Teaching Effectiveness），本研究从 \texttt{run\_01} 中随机抽取了多段对话记录（\texttt{manual\_audit.csv}）进行人工与大语言模型裁判（LLM-as-a-Judge）的双重评估。评估维度主要聚焦于“是否成功引导学生产生认知顿悟”。

在对 \texttt{sim\_Baseline\_P1\_M-ELE-001\_8db7b7} 的审核中，Baseline 系统（评分为 2 分）虽然频繁使用反问和类比（如水车、送货员等），但始终未能切中学生“电能与电荷流动”混淆的要害。由于缺乏阶段控制，系统在学生仍坚持错误观念时便陷入了比喻的堆砌，最终未能帮助 P1（固执型）学生突破认知僵局。

相比之下，在 \texttt{sim\_FSM+Guardrail\_P2\_M-BUO-001\_809f87}（浮力重物必沉迷思）中，系统精准地执行了阶段化控制：在 S4 阶段通过 Assumption\_Probing 将“重的必沉”这一经验规则显化，并以极端反例（“钢铁巨轮为何能浮，小铁钉为何沉底？”）制造认知冲突；随后在 S5/S6 阶段把讨论从“重量”迁移到“体积$\rightarrow$排水$\rightarrow$托力”，用受控思想实验支架（“如果把巨轮揉成实心铁球会怎样？”）帮助学生主动说出关键机制。该轨迹在人工抽核中获得了较高的教学有效性评分，充分体现了阶段化控制在复杂概念教学中的优势。

\section{教学安全约束机制的效果分析}
\label{sec:5.3}

\subsection{答案泄露率（Answer Leakage Rate）对比}
\label{subsec:5.3.1}

生成式人工智能在教育应用中的核心风险之一是沦为“作业代写机”，直接向学生泄露最终答案，从而剥夺学生的认知劳动。本研究在第四章中设计了输出侧认知卸载审查机制。表~\ref{tab:leakage-rate} 呈现了该机制拦截答案泄露的直接效果。

\begin{table}[htbp]
    \centering
    \caption{各版本系统答案泄露率对比}
    \label{tab:leakage-rate}
    \begin{tabular}{lccc}
        \hline
        \textbf{版本} & \textbf{答案泄露率均值} & \textbf{标准差} & \textbf{95\% 置信区间} \\
        \hline
        Baseline & 7.13\% & 4.29\% & $\pm 3.76\%$ \\
        FSM & 8.94\% & 1.77\% & $\pm 1.55\%$ \\
        FSM+Guardrail & 0.00\% & 0.00\% & $\pm 0.00\%$ \\
        \hline
    \end{tabular}
\end{table}

数据清晰地表明，未加护栏约束的 Baseline 和单纯状态机控制的 FSM 版本，其答案泄露率分别达到了 7.13\% 和 8.94\%。在这些版本中，当学生反复表达困惑或直接索要公式时，大语言模型往往会出于其内置的“帮助人类用户完成任务”的对齐倾向，直接输出完整的推导过程或最终结论。

而引入安全约束机制的 FSM+Guardrail 版本，其答案泄露率被完美控制在 0.00\%。这一结果强有力地证明了我们在第四章中设计的输出审查漏斗（基于 LLM-as-a-Judge 的 Leakage Evaluation）能够严格守住“不直接给答”的底线。

\subsection{拦截触发特征与启发式豁免效果分析}
\label{subsec:5.3.2}

“零泄露”并非通过简单粗暴的“拒绝回答”来实现的。根据聚合数据，FSM+Guardrail 版本的平均护栏拦截触发率（Guardrail Interception Rate）为 2.60\%（$\pm 1.35\%$）。这意味着系统并非时刻在拦截，而是在关键节点上发挥了兜底作用。

更重要的是，第四章中提出的“教学支架豁免”机制在保持低泄露率的同时，维系了对话的流畅度。在 \texttt{sim\_FSM+Guardrail\_P3\_M-BUO-002\_aade88}（压强经验投射迷思）日志中，面对 P3（易挫败型）学生在上下表面压力差问题上的卡顿，系统成功触发了【解释类比/思想实验】的豁免：“假设原来上面受压 10N，下面受压 15N...现在都增加 100N...新的差值是多少？”。系统通过给出具体的计算支架，帮助学生跨越了抽象符号推演的鸿沟，而没有直接告诉其“浮力等于排开液体重力”的最终结论。这种在“安全边界内提供最小充分支架”的策略，有效缓解了认知过载，避免了对话因过度严格的拦截而陷入死锁。

\subsection{自动化模拟评估与人工抽核：苏格拉底式提问特征符合度对比}
\label{subsec:5.3.3}

在对微观干预日志的审计中（\texttt{manual\_audit.csv}），大语言模型裁判对所有版本的“苏格拉底度（Socratic Degree）”均给出了最高分（5分）。这表明无论是否有护栏，基础模型在提示词工程的引导下，都具备了采用反问、类比和启发式探究的基本语域能力。

然而，护栏版本的优越性体现在提问的“韧性”上。例如，在 \texttt{sim\_FSM\_P3\_M-ELE-001\_fd94e4} 案例中，纯 FSM 版本虽然也是以提问为主，但当对话陷入深水区时，出现了多达 4 次的轻微答案泄露。而在对应的 FSM+Guardrail 版本中，系统不仅保持了苏格拉底式的循循善诱，还在学生面临“顿悟”的临界点上，坚守住了底线，逼迫学生自行说出关键机制（如“原来只要托力够大，就算很重也能浮着啊”），从而实现了更高质量的认知参与。

\section{对话干预的差异表现与边界现象}
\label{sec:5.4}

\subsection{阶段化控制+护栏版本下的优质干预案例}
\label{subsec:5.4.1}

本节在宏观均值之外，从 \texttt{session\_summary.jsonl} 与 \texttt{turn\_logs.jsonl} 进一步给出按主题（电学/浮力）与学生画像（P1/P2/P3）的分组统计，以识别“总体均值被掩盖”的结构性差异。

\begin{table}[htbp]
    \centering
    \caption{按主题与画像分组的认知修正率（合并五次运行；括号内为会话数）}
    \label{tab:grouped-correction-topic-profile}
    \begin{tabular}{llccc}
        \hline
        \textbf{版本} & \textbf{主题} & \textbf{P1} & \textbf{P2} & \textbf{P3} \\
        \hline
        Baseline & 电学 & $0.0\%~(n=10)$ & $70.0\%~(n=10)$ & $80.0\%~(n=10)$ \\
        Baseline & 浮力 & $0.0\%~(n=10)$ & $100.0\%~(n=10)$ & $90.0\%~(n=10)$ \\
        \hline
        FSM & 电学 & $0.0\%~(n=10)$ & $100.0\%~(n=10)$ & $90.0\%~(n=10)$ \\
        FSM & 浮力 & $0.0\%~(n=10)$ & $100.0\%~(n=10)$ & $70.0\%~(n=10)$ \\
        \hline
        FSM+Guardrail & 电学 & $0.0\%~(n=10)$ & $80.0\%~(n=10)$ & $90.0\%~(n=10)$ \\
        FSM+Guardrail & 浮力 & $0.0\%~(n=10)$ & $100.0\%~(n=10)$ & $80.0\%~(n=10)$ \\
        \hline
    \end{tabular}
\end{table}

\begin{figure}[htbp]
    \centering
    \fbox{\parbox{0.92\linewidth}{\centering 图待补绘：按主题（电学/浮力）分面展示\\
    不同版本在 P1/P2/P3 上的认知修正率柱状图（含 95\% CI 或自助法误差线）。}}
    \caption{按主题与画像分组的认知修正率差异（图待补绘）}
    \label{fig:grouped-correction-topic-profile}
\end{figure}

绘图要点：横轴为画像（P1/P2/P3），纵轴为认知修正率；用颜色区分版本（Baseline/FSM/FSM+Guardrail），并按主题（电学/浮力）做分面或子图；误差线建议以 Run 为单位做自助法重采样或直接使用各 Run 的比例计算 95\% CI。

表~\ref{tab:grouped-correction-topic-profile} 显示出两个对总体结论具有解释力的现象。第一，P1 在两主题与三版本下的修正率均为 0\%，说明“固执型/逻辑熔断”画像对纯对话式干预构成了主导性难点，其失败并非某个具体机制缺失即可单点修复，而更接近“交互形态与权限上限”导致的结构性瓶颈。第二，FSM 对 P2/P3 的提升在电学主题更为显著：电学任务更依赖“闭合回路/能量与电荷区分”等层级化概念链条，阶段化推进（T1--T5）能够更稳定地组织冲突、支架与验证序列；而在浮力主题中，部分反例与数值类比对 Baseline 也足够有效，因此 FSM 的边际收益相对收敛。对 FSM+Guardrail 而言，其主要增益体现在“零泄露”而非“更高修正率”，这与前述安全约束与教学效率之间的 trade-off 一致（见表~\ref{tab:correction-rate} 与表~\ref{tab:leakage-rate}）。

\subsection{成功案例：FSM+Guardrail 在 P2 浮力主题的认知冲突—支架链路}
\label{subsec:5.4.2}

本节选取 \texttt{sim\_FSM+Guardrail\_P2\_M-BUO-001\_809f87}（见 \texttt{run\_01/results/manual\_audit.csv} 与对应 \texttt{turn\_logs.jsonl}）作为成功案例，对其“机制$\rightarrow$过程$\rightarrow$结果”的证据链进行补强。

\paragraph{选取理由}该会话在人工抽核中被标注为高教学有效性（有效性 5 分），且完整覆盖了从错误规则暴露（“重的必沉”）到反例冲突、再到以学生自述形式产出关键判别关系（$F_{\text{浮}}=G$）的闭环；同时其全程无护栏拦截与无泄露，能较纯粹地体现阶段化控制与策略路由对“认知推进”的贡献。

\paragraph{关键轮次}第 1 轮系统以“巨轮 vs 铁钉”构造极端反例，迫使学生承认经验规则无法解释现实（S4，Assumption\_Probing）。第 3--4 轮进一步将讨论从“重量”转移到“形状$\rightarrow$体积$\rightarrow$排水量”，并用“揉成实心铁球必沉”的反事实推演稳定冲突张力（S6，Evidence\_Seeking/Consequence\_Exploration）。第 5--6 轮系统引导学生自主说出“排开水越多托力越大”与“漂浮时二力平衡”，并追问铁钉为何无法获得足够托力以完成迁移。

\paragraph{机制对应}该过程对应第四章的阶段化教学主线：在 T3（认知冲突）阶段使用后果探索/反事实推演策略稳固矛盾，再在 T4（支架过渡）阶段将“重量”这一表层变量重构为“排开体积$\rightarrow$浮力”这一结构变量，并通过连续证据追问（Evidence\_Seeking）把关键关系以学生口头表达的方式外显，从而避免系统直接讲授结论造成的认知卸载。

\paragraph{结果}该会话共 6 轮即达成 \texttt{resolved} 终止，\texttt{resolved\_flag=true} 且 \texttt{answer\_leakage\_count=0}、\texttt{guardrail\_trigger\_count=0}；其轨迹体现出“用最短路径抵达可验证理解”的理想形态，也解释了表~\ref{tab:system-stability} 中 FSM+Guardrail 平均轮次更短且波动更小的现象。

\subsection{边界/失败案例：P1 电学主题下的逻辑熔断与文本对话上限}
\label{subsec:5.4.3}

尽管系统在多数情况下表现稳定，但针对 P1（固执型/逻辑熔断易发者）的干预呈现出系统性边界：在本次套件的两主题与三版本下，P1 的认知修正率均为 0\%（表~\ref{tab:grouped-correction-topic-profile}）。以下以 \texttt{sim\_FSM+Guardrail\_P1\_M-ELE-002\_c91ae9}（见 \texttt{run\_05/results/manual\_audit.csv} 与对应 \texttt{turn\_logs.jsonl}）为例说明其失败机制。

\paragraph{选取理由}该会话满足“失败但非系统错误”的典型性：全程无异常终止、无护栏拦截、无答案泄露，却在 13 轮达到上限并以 \texttt{max\_turns\_reached} 结束，说明瓶颈主要来自“学生模型的抵抗性”与“文本交互的证据供给不足”，而非工程稳定性或安全策略误伤。

\paragraph{关键轮次}第 1--3 轮系统尝试用反事实与可观察现象（断开回路灯泡不亮）推动学生承认“仅接正极”不可行（S4$\rightarrow$S5），但学生将失败归因于“绝缘/接触不良”，并坚持“电是能量、用完就消失”。第 5--6 轮系统引入拥堵类比（地铁单向进站）试图显化“载体堆积/无法持续流动”的矛盾，学生明确拒绝类比迁移（“电不是人”）。第 7 轮系统转而采用事实落地（Fact\_Grounding），强调实验结论并请求学生类比吸管堵塞体验，但学生仍以“电不同于空气”回避。此后第 9--13 轮系统在水泵类比、思想实验与追问之间循环，学生回复持续回到“能量被消耗$\rightarrow$无需回路”的单点解释，未形成可被验证阶段吸收的新表征。

\paragraph{机制对应}该失败链条对应第四章中“阶段化控制在高抵抗性画像下可能触发逻辑熔断”的预期：当学生对类比与反事实均不接受时，系统缺乏多模态实证（可操作电路、动态仿真）与更高权限讲授手段，只能在 T3/T4 的语言性支架内反复加压；护栏机制虽成功守住“不直接给答”的底线，但也意味着系统无法用“直接讲授”作为最后手段完成概念覆盖，从而使死锁更容易以轮次耗尽的形式显现。

\paragraph{结果}该会话最终 \texttt{resolved\_flag=false}，后测未通过，终止原因为 \texttt{max\_turns\_reached}。该类失败并不否定阶段化对话与护栏机制的总体价值，而是提示：在 P1 等高抵抗性场景下，仅文本的苏格拉底式提问可能不足以触发概念替换，需要引入更强的证据形态或更明确的教学权限切换，以满足概念转变理论所强调的“可理解性、可接受性与生成性”条件\cite{posner1982accommodation,chi1994things}。

\section{有效性威胁、局限与改进方向}
\label{sec:5.5}

本章结果建立在一次固定测试套件的五次重复运行之上，仍存在以下有效性威胁与局限。
\begin{itemize}
    \item \textbf{样本结构与外推性}：本套件的“主题$\times$画像$\times$迷思”组合是对真实课堂的抽象与离散化模拟，尤其是 P1/P2/P3 的行为生成规则可能与真实学生群体存在系统偏差。因此本章更强调“机制在控制条件下是否按预期工作”，而非直接宣称可对任意真实教学场景外推其绝对增益。
    \item \textbf{LLM-as-a-Judge 偏差与一致性风险}：\texttt{evaluation\_results.json} 的苏格拉底度与有效性评分依赖裁判模型，其标注偏好可能与人类教师评价标准不完全一致。为降低单一裁判带来的系统偏差，本章在关键断言处优先使用可复核的行为指标（如 \texttt{answer\_leakage\_flag}、\texttt{guardrail\_triggered}、\texttt{resolved\_flag}），并保留 \texttt{manual\_audit.csv} 作为人工抽核入口。
    \item \textbf{指标间 trade-off 的解释边界}：护栏带来的“零泄露”与教学推进效率之间存在可观的权衡空间。由于本套件并未系统扫参（例如不同护栏阈值、不同重生成次数），因此无法在统计意义上给出最优策略的充分比较，只能就当前实现阐明其行为后果。
    \item \textbf{文本对话的证据供给上限}：如第~\ref{subsec:5.4.3} 节所示，当学生拒绝类比迁移、拒绝反事实推演或持续诉诸单点经验规则时，纯文本互动缺少可操作实验与可视化反馈，难以满足概念转变对“冲突证据强度”的需求。这提示未来迭代应优先考虑与项目路线一致的能力增强：在不扩大越权给答风险的前提下，引入可控的多模态演示（虚拟电路/浮力仿真）与更精细的“权限升级”策略（例如在连续死锁后切换为更明确的讲授模板，同时保持可追溯的安全审查）。 
\end{itemize}
